% =============================================================================
%  Related Work — SUDA-Muon
%  Target venues: ICML / NeurIPS / ICLR
%  Sections:
%    1. Muon and Matrix-Aware Optimizers
%    2. Decentralized Optimization & Heterogeneity Correction
%    3. Decentralized Muon (The Gap)
% =============================================================================

\section{Related Work}
\label{sec:related_work}

% ---------------------------------------------------------------------------
\subsection{Muon and Matrix-Aware Optimizers}
\label{subsec:muon}

Modern deep learning architectures—Transformers, MLPs, and their variants—are
parameterized almost entirely by \emph{matrix-structured} weight tensors.
Classical first-order methods such as SGD~\citep{nedic2009distributed} and
Adam~\citep{kingma2014adam} treat every parameter entry symmetrically, ignoring
the intrinsic operator geometry of these matrices.  A growing body of work
argues that \emph{matrix-aware} update rules, which respect the spectral
structure of weight matrices, yield substantially faster convergence and better
generalization.

\paragraph{Muon.}
The Muon optimizer~\citep{jordan2024muon} replaces the raw stochastic gradient
$G$ with its \emph{polar factor} (matrix sign)
$\operatorname{msgn}(G) = UV^{\top}$, where $G = U\Sigma V^{\top}$ is the
singular value decomposition.  This \emph{scale-invariant} update aligns the
descent direction with the geometry of the operator norm rather than the
Euclidean norm.  In practice, the polar factor is approximated via a small
number of Newton--Schulz iterations, avoiding the cost of a full
SVD~\citep{amsel2025polar, grishina2025chebyshev}.
\citet{liu2025muon_scalable} demonstrate that Muon, augmented with weight decay
and per-layer learning-rate adjustment, scales to billion-parameter language
model pre-training, matching or surpassing AdamW across a wide range of model
sizes.

\paragraph{Theoretical foundations.}
Several recent works have placed Muon on firm theoretical ground.
\citet{li2025note} show that Muon converges at the rate
$\mathcal{O}(1/\sqrt{T})$ for non-convex stochastic objectives under standard
assumptions, establishing it as a provably valid optimizer.
\citet{chen2025muon_spectral} prove that Muon implicitly solves a
spectral-norm-constrained trust-region subproblem at each step, providing a
clean geometric interpretation.
\citet{kovalev2025gradient} offer a complementary analysis via non-Euclidean
trust-region theory, showing that orthogonalized gradient steps are optimal
under an operator-norm ball constraint.
\citet{polargrad2025} unify Muon and related methods (Shampoo, SOAP) under the
\emph{PolarGrad} preconditioning framework, elucidating the shared algebraic
structure.

\paragraph{Extensions and variants.}
The Gluon optimizer~\citep{riabinin2025gluon} bridges the gap between the
theoretical LMO (Linear Minimization Oracle) perspective and practical
implementation by deriving corrected step-size schedules for Muon-type methods.
Dion~\citep{ahn2025dion} extends orthonormalized updates to the sharded-weight
regime of large-scale distributed training, reducing the communication overhead
of dense matrix operations.
\citet{pethick2025norm} study a broader family of norm-constrained LMO-based
optimizers, characterizing their convergence under unified assumptions.
Despite this rapid progress, \emph{all} existing theoretical analyses assume a
\textbf{centralized} or \textbf{single-node} setting.  The question of how
Muon's non-linear polarization interacts with a peer-to-peer communication
topology is left open.

% ---------------------------------------------------------------------------
\subsection{Decentralized Optimization and Heterogeneity Correction}
\label{subsec:decentralized}

\paragraph{Decentralized SGD.}
Decentralized Parallel SGD (D-PSGD)~\citep{lian2017can} avoids the
communication bottleneck of a central parameter server by having each agent
average its iterate with neighbors over a fixed mixing matrix.
\citet{lian2017can} show that D-PSGD achieves \emph{linear speedup} in the
number of agents $N$: the dominant $\mathcal{O}(1/\sqrt{NT})$ term matches
centralized mini-batch SGD, while the network-topology penalty appears only in
lower-order terms.  This result has been generalized to time-varying graphs,
partial participation, and other practical settings~\citep{koloskova2020unified}.
The practical importance of decentralized training at scale is underscored by
work on training foundation models without a central
server~\citep{yuan2022decentralized}.

\paragraph{The client-drift problem.}
When agents hold heterogeneous (Non-IID) data, pure D-PSGD suffers from
\emph{client drift}: local gradients are biased toward each agent's own
distribution, causing iterates to diverge from the global optimum.
\citet{tang2018d2} quantify this bias in the $\mathrm{D}^2$ algorithm and show
that it can dominate the convergence error.
\citet{yuan2021removing} demonstrate that data heterogeneity effectively
inflates the apparent spectral gap of the mixing matrix, masking the true
topology dependence of D-SGD.

\paragraph{Heterogeneity-correction methods.}
A rich literature has developed exact, bias-free decentralized algorithms.
EXTRA~\citep{shi2015extra} introduces a \emph{difference of mixing matrices}
correction that cancels the gradient bias at fixed step-sizes, achieving exact
convergence for smooth convex objectives.
Exact Diffusion / $\mathrm{D}^2$~\citep{tang2018d2} reinterprets this
correction through an augmented Lagrangian lens, connecting decentralized
methods to primal-dual splitting.
Gradient Tracking (GT)~\citep{pu2020distributed} maintains a local estimate of
the global gradient average and updates it in tandem with the primal variable;
\citet{xin2020improved} establish $\mathcal{O}(1/\sqrt{NT})$ non-convex
convergence for stochastic GT.

\paragraph{Unified primal-dual framework (SUDA).}
\citet{alghunaim2022unified} show that EXTRA, Exact Diffusion, and Gradient
Tracking are all special instances of a single \emph{primal-dual recursion}
parameterized by low-degree polynomials of the mixing matrix.  This SUDA
(Stochastic Unified Decentralized Algorithm) framework yields a \emph{refined}
convergence analysis in which data heterogeneity affects only higher-order
terms, so the leading $\mathcal{O}(1/\sqrt{NT})$ rate is topology-independent
and matches centralized SGD.  Critically, SUDA is designed for \emph{Euclidean}
gradient updates; accommodating the non-linear matrix orthogonalization
required by Muon demands a non-trivial extension.

\paragraph{Decentralized Riemannian optimization.}
Decentralized optimization over the Stiefel manifold has been studied
by~\citet{chen2021decentralized}, who combine Riemannian gradient steps with
gossip averaging and establish convergence for smooth non-convex objectives.
While Muon's update can be viewed as a retraction onto the Stiefel manifold
(via the polar factor), the stochastic, large-scale regime of LLM training
and the need for a unified multi-scheme analysis go beyond the scope of
existing Riemannian decentralized methods.

% ---------------------------------------------------------------------------
\subsection{Decentralized Muon and the Motivating Gaps}
\label{subsec:demuon}

\paragraph{DeMuon: the pioneering work.}
DeMuon~\citep{he2025demuon} is the first algorithm to extend Muon to a
decentralized graph.  It combines \emph{gradient tracking} with Newton--Schulz
orthogonalization: each agent tracks the global gradient average, applies the
matrix sign operator locally, and then mixes iterates with neighbors.
DeMuon establishes convergence for non-convex stochastic objectives and
demonstrates empirical gains over decentralized Adam on matrix-structured tasks.

\paragraph{Gap 1 — Algorithmic rigidity.}
DeMuon is tied to a \emph{single} communication template: gradient tracking.
It cannot, without redesign, incorporate other powerful heterogeneity-correction
schemes such as EXTRA or Exact Diffusion.  The SUDA framework shows that
different polynomial choices of the mixing matrix can yield superior practical
performance on specific network topologies; a Muon variant that is locked to
one choice foregoes these benefits.  Our work introduces a \emph{unified}
primal-dual perspective that accommodates any SUDA-compatible communication
template.

\paragraph{Gap 2 — Topology-separated analysis.}
Existing analyses for decentralized Muon (and for DeMuon in particular) do not
exploit the refined, topology-separated convergence structure established by
SUDA-style frameworks.  Concretely, the leading convergence term in DeMuon's
analysis mixes the spectral gap of the network with the data heterogeneity
constant, preventing a clean separation between optimization difficulty and
network quality.  SUDA-Muon, by contrast, recovers a rate where the network
penalty appears \emph{only} in higher-order terms, matching the topology
sensitivity of centralized Muon.

\paragraph{Gap 3 — Linear speedup.}
Linear speedup—the property that adding $N$ agents reduces the dominant error
term by a factor of $N$—is a cornerstone guarantee of decentralized SGD
\citep{lian2017can, koloskova2020unified}.  It has also been established for
centralized Muon~\citep{li2025note}.  However, whether a \emph{serverless}
Muon algorithm can achieve asymptotic linear speedup in $N$ remained unknown
prior to our work.  The non-linearity of the polar-factor operator complicates
the standard noise-averaging argument, and DeMuon's analysis does not resolve
this question.  SUDA-Muon is the first decentralized Muon algorithm with a
provable $\mathcal{O}(1/\sqrt{NT})$ linear-speedup guarantee.

\paragraph{Related distributed matrix optimizers.}
Concurrent work on distributed orthonormalized updates includes
Dion~\citep{ahn2025dion}, which focuses on the data-parallel (centralized
server) setting with sharded weights, and the low-rank orthogonalization
framework of~\citet{he2025lowrank}, which targets foundation-model fine-tuning
rather than serverless decentralized training.  Neither provides the
topology-separated convergence theory or the multi-scheme unification that
motivates SUDA-Muon.
